% Title, one-page summary, reading guide and table of contents.
\thispagestyle{plain}
\begin{center}
{\color{accentdark}\rule{\linewidth}{0.8pt}}\\[10pt]
{\LARGE\bfseries Repo implicite Totem}\\[6pt]
{\large Contrôler le repo du front office et détecter les courbes aberrantes}\\[8pt]
{\normalsize\color{ink2} Note méthodologique \quad\textbullet\quad Contrôle indépendant des prix, dérivés
actions}\\[2pt]
{\normalsize\color{ink2} Octobre 2026 \quad\textbullet\quad Document de travail}\\[4pt]
{\color{accentdark}\rule{\linewidth}{0.8pt}}
\end{center}

\begin{bref}
\textbf{Le problème.} Le contrôle indépendant des prix déduit de Totem une courbe de repo implicite : les forwards
viennent des prix de consensus des calls et des puts, les dividendes d'un service Totem distinct. Les deux services
n'ont pas les mêmes contributeurs, et leur nombre varie avec la maturité ; dans notre exemple, \vRunNoptShort{}
contributeurs options et \vRunNdivShort{} contributeurs dividendes à court terme, \vRunNoptLong{} et \vRunNdivLong{} à
cinq ans. La courbe de repo implicite hérite du bruit des deux panels : elle paraît parfois irrégulière, ou saute
d'une maturité à l'autre sans raison de marché. C'est pourtant elle qui sert, par l'intermédiaire des
sensibilités, à contrôler les positions repo du front office (FO).

\medskip
\textbf{Les trois objectifs de la note.}
\begin{description}[leftmargin=1.1cm,labelwidth=0.9cm,labelsep=0.2cm,itemsep=1pt]
\item[O1] \emph{Comprendre} : d'où viennent ces formes, et lesquelles ne sont que du bruit normal ?
\item[O2] \emph{Contrôler} : \og j'ai le contrôle forward, j'ai le contrôle dividendes : combien vaut le contrôle
repo ? \fg{} Peut-on contrôler la position repo sans dépendre des contributeurs dividendes ?
\item[O3] \emph{Détecter} : comment signaler automatiquement une courbe Totem aberrante, en nommer la cause, et ne
pas la confondre avec une erreur de marquage du FO ?
\end{description}

\medskip
\textbf{Les réponses.}
\begin{description}[leftmargin=1.1cm,labelwidth=0.9cm,labelsep=0.2cm,itemsep=2pt]
\item[R1] Le repo implicite est égal au repo moyen du panel options, plus un \emph{écart de panel} : la différence
entre deux moyennes de dividendes calculées sur des contributeurs différents. Pour une action européenne,
cet écart vaut couramment \vGapOneY\,pb à un an et \vGapFiveY\,pb à cinq ans (\cref{sec-diagnostic}). Aucune méthode qui ne voit
que les données Totem ne peut le retirer ; on peut en revanche en calculer la taille et reconnaître, à leur forme,
les mécanismes qui produisent les vraies anomalies.
\item[R2] Ligne repo $=$ ligne forward $-$ ligne dividendes $+$ ligne non-forward. Pour un portefeuille dont la
valeur ne dépend que des forwards (options européennes, forwards, futures), la ligne non-forward est nulle : le
contrôle total ne dépend pas des contributeurs dividendes, et il en est exactement indépendant en réévaluation
complète (\cref{sec-controle}). Un vecteur d'exposition calculé à partir des seules sensibilités FO indique où la
répartition entre repo et dividendes compte.
\item[R3] Le détecteur structurel juge l'écart Totem--FO \og en portage \fg{}, avec une bande de bruit calculée à
partir des dispersions et des nombres de contributeurs publiés par Totem, et teste une forme type par mécanisme
(\cref{sec-detection}). Son taux de fausses alertes est fixé (\vFPR{} des sous-jacents sains) ; chaque alerte est
expliquée en unités naturelles et accompagnée d'une courbe réparée. Sur un banc d'essai synthétique de \vNTest{}
sous-jacents (\cref{sec-banc}), il détecte en moyenne autant que la meilleure alternative testée, B-spline et mélange
gaussien (\vMeanStruct{} contre \vMeanGmmDiff{}), trouve bien mieux les forwards aberrants (\vFwdStruct{} contre
\vFwdGmmDiff{}) et ne prend pas une erreur de marquage FO pour une courbe Totem aberrante (\vFoStruct{} d'alertes
contre \vFoGmmDiff{}).
\end{description}
\end{bref}

\newpage
\section*{Comment lire cette note}
Chaque section commence par son objectif et se termine par un encadré \og À retenir \fg{}. Un même sous-jacent
synthétique, l'\emph{exemple fil rouge}, sert d'illustration du début à la fin (encadrés verts) : une action européenne
observée d'abord dans un marché sain, puis sous huit anomalies différentes. Les démonstrations courtes suivent les
énoncés ; les plus longues sont regroupées en \cref{app-proofs}.

\begin{center}\small
\begin{tabular}{>{\raggedright\arraybackslash}p{5.1cm}c>{\raggedright\arraybackslash}p{8.0cm}}
\toprule
Question & Section & Ce que l'on y trouve\\
\midrule
Que cherche-t-on, et pourquoi est-ce difficile ? & \ref{sec-probleme} & La chaîne de contrôle, l'effet loupe en $1/T$,
l'exemple fil rouge\\
Quelles notations ? & \ref{sec-cadre} & Forward, écart de portage, données Totem, passage au repo implicite\\
\textbf{O1} D'où viennent les courbes aberrantes ? & \ref{sec-diagnostic} & L'écart de panel et sa taille, les sorties
de contributeurs, six mécanismes et leurs signatures, ce que les données ne peuvent pas dire\\
\textbf{O2} Comment contrôler le repo ? & \ref{sec-controle} & L'identité de contrôle, trois portefeuilles,
la matérialité de la répartition repo/dividendes\\
\textbf{O3} Comment détecter ? & \ref{sec-detection} & Le modèle de bruit, les formes types, le test, la calibration,
le diagnostic de l'exemple fil rouge\\
Est-ce que ça marche ? Et les B-splines, les mélanges, les réseaux de neurones ? & \ref{sec-banc} & Le banc d'essai
et la comparaison avec les idées de départ\\
Que faire en pratique ? & \ref{sec-reco} & Recommandations, données nécessaires, procédure mensuelle\\
\bottomrule
\end{tabular}
\end{center}

{\small\setcounter{tocdepth}{2}\tableofcontents}
\newpage
